\chapter{Metodología}

\section{Enfoque y tipo de investigación}

El trabajo corresponde a una investigación aplicada de desarrollo tecnológico con
validación experimental en dispositivo. No se busca generar conocimiento teórico
nuevo sobre detección de URLs maliciosas, materia ya tratada en los antecedentes,
sino construir un artefacto funcional que lleve capacidades de análisis existentes
al único punto de la cadena del ataque donde todavía se pueden aplicar: el intervalo
entre la lectura del código QR y la apertura del navegador. El método, por lo tanto,
se organiza alrededor de dos preguntas verificables: si el artefacto cumple los
tiempos y las tasas de acierto que lo vuelven usable, y si el veredicto que produce
se entiende lo suficiente como para que el usuario decida bien.

La unidad de análisis es exclusivamente la URL textual extraída del código, según el
RFC~3986, tal como se delimitó en el capítulo de alcance. No se analiza el símbolo QR
como imagen ni ningún otro contenido codificable.

Conviene fijar desde el inicio una distinción que ordena todo el diseño de la
evaluación. El presente trabajo mide el desempeño del módulo de frontend: la latencia
del escaneo, la tasa de detección de códigos, la corrección de la clasificación de
contenido, la integridad del flujo y la comprensión del veredicto por parte del
usuario. No mide la capacidad de detección del motor de análisis, que depende de las
fuentes de reputación y de los filtros implementados en el módulo de backend. Son
dos conjuntos de criterios distintos y ninguno sustituye al otro: una pantalla muy
clara que comunique un veredicto poco confiable no protege a nadie, y un veredicto
acertado pero mal explicado tampoco. Cuando una métrica dependa del backend, se
declara explícitamente en la sección correspondiente.

\section{Arquitectura de la solución}

El prototipo se organiza en cuatro capas, cada una con una responsabilidad que no se
solapa con las demás. Esta separación no es ornamental: es la que permite que la
lógica de negocio se pruebe sin cámara, sin red y sin dispositivo, lo que a su vez
determina qué parte de la validación puede automatizarse y qué parte exige hardware.

\begin{table}[H]
    \centering
    \caption{Capas de la aplicación y su responsabilidad}
    \begin{tabular}{@{}>{\raggedright\arraybackslash}p{3.2cm}>{\raggedright\arraybackslash}p{8.5cm}@{}}
        \toprule
        \textbf{Capa} & \textbf{Responsabilidad} \\
        \midrule
        Rutas & Doce pantallas con enrutamiento basado en archivos: escaneo, carga, resultado, detalle, contenido no navegable, historial, reporte, Story Card y dos pasos de onboarding \\
        Sistema de diseño & Diecisiete componentes y los tokens de color, tipografía y espaciado, extraídos del prototipo visual de alta fidelidad \\
        Servicios & Lógica de negocio pura y aislada: análisis, cliente HTTP, clasificación de QR, almacenamiento, ubicación, reporte y formato \\
        Estado & Un único contexto de React que sostiene el resultado vigente, el historial y la preferencia de ubicación \\
        \bottomrule
    \end{tabular}
    \label{tab:capas}
\end{table}

La dirección de las dependencias merece una nota, porque es contraria a lo que suele
esperarse. El tipo que enumera los cuatro veredictos posibles ---seguro, sospechoso,
malicioso e incierto--- reside en el sistema de diseño, dado que cada veredicto lleva
asociados color, etiqueta y tratamiento visual. En consecuencia, la capa de servicios
depende del sistema de diseño y no al revés. La implicación práctica es que un quinto
veredicto se agregaría en el sistema de diseño y desde ahí se propagaría al resto.

\subsection{Plataforma y versiones}

El prototipo se construye sobre el stack que recoge el cuadro~\ref{tab:stack-metodo}.
Se consignan las versiones exactas porque varias decisiones del método, en particular
la que se discute en la sección siguiente, se explican por restricciones de
compatibilidad entre versiones concretas y no por preferencias generales.

\begin{table}[H]
    \centering
    \caption{Versiones verificadas del entorno de desarrollo}
    \begin{tabular}{@{}>{\raggedright\arraybackslash}p{5.2cm}>{\raggedright\arraybackslash}p{4.0cm}@{}}
        \toprule
        \textbf{Componente} & \textbf{Versión} \\
        \midrule
        Expo SDK & 57.0.12 \\
        React Native & 0.86.2 (New Architecture) \\
        React & 19.2.3 \\
        Enrutamiento por archivos & Expo Router 57.0.12 \\
        Cámara y escaneo & expo-camera 57.0.4 \\
        Almacenamiento clave-valor & expo-sqlite 57.0.2 \\
        Lenguaje & TypeScript 6.0.3, modo \texttt{strict} \\
        Pruebas & Jest 30.5.1 con preset \texttt{jest-expo} \\
        Entorno de compilación & Node 20.20.2, JDK 17, Gradle 9.3.1, API mínima 24 \\
        \bottomrule
    \end{tabular}
    \label{tab:stack-metodo}
\end{table}

\subsection{Decisión sobre las librerías de escaneo y almacenamiento}

El anteproyecto especificaba \texttt{react-na\-ti\-ve-vision-camera} para el escaneo y
MMKV para el almacenamiento local. La implementación emplea los módulos nativos del
propio SDK: \texttt{expo-camera} y el almacén clave-valor de \texttt{expo-sqlite}. La
desviación se documenta aquí con su justificación técnica y con el criterio que la
revertiría, porque un cambio de esta naturaleza debe quedar declarado como decisión y
no descubrirse al leer el código.

El motivo es un conflicto de dependencias nativas verificado durante la fase de
integración. El enrutador por archivos arrastra la versión 0.12.x de la librería de
\emph{worklets}, mientras que el núcleo de módulos del SDK en su versión 57.0.10 solo
admite versiones iguales o anteriores a la 0.10.x. La combinación rompe la
compilación nativa de Android con un error de símbolo ausente en tiempo de enlace, lo
que obligó a fijar versiones explícitas de resolución en el manifiesto de
dependencias. Adoptar los módulos del propio SDK reduce esa superficie de fallo, pues
sus versiones se publican alineadas con el SDK, sin sacrificar ninguna de las
capacidades que el escaneo requiere: decodificación nativa en tiempo real restringida
al formato QR, control de linterna y gestión de permisos de cámara.

En el caso del almacenamiento, la razón es funcional además de práctica. El almacén
clave-valor de \texttt{expo-sqlite} ofrece una interfaz \emph{síncrona} de lectura.
Esto permite resolver dos decisiones durante el primer renderizado: a qué pantalla
dirigir al usuario según si ya completó el onboarding, y con qué contenido pintar el
historial. Con una interfaz asíncrona, ambas decisiones ocurrirían después del primer
renderizado y producirían un parpadeo visible hacia la pantalla equivocada.

El criterio de reversión está fijado de antemano y es el resultado de la prueba~1.4,
descrita en la sección~\ref{sec:mediciones}. Si la tasa de detección resulta inferior
al 95~\% o el percentil~90 del tiempo de escaneo supera los 5~s, se migra a
\texttt{react-native-vision-camera} conforme al anteproyecto original, con un esfuerzo
estimado de uno a dos días. Antes de concluir que la migración es necesaria se
descartan tres hipótesis de menor costo: que el fallo se concentre en ángulos de
45\textdegree, caso en que correspondería ampliar el área de escaneo; que se concentre
en penumbra, caso en que la linterna podría resolverlo y la interfaz debería
sugerirla; o que se concentre en códigos pequeños a distancia, que sería una
limitación de enfoque y no de decodificación. Si el fallo se concentra en un solo eje
se documenta como limitación conocida; si se reparte en los tres, la migración queda
justificada.

\section{Procedimiento de desarrollo}

El desarrollo se organiza en ocho bloques ordenados por dependencia técnica y riesgo,
no por la secuencia de fases del anteproyecto. El cambio de criterio responde a una
situación concreta: al iniciar la integración existían tareas de fases avanzadas
terminadas y tareas de fases iniciales todavía abiertas, de modo que el orden original
había dejado de describir el estado real del proyecto. La reorganización se rige por
tres reglas declaradas: primero lo que desbloquea el trabajo de otras personas,
después lo que elimina la mayor cantidad de riesgo técnico y al final lo que depende
únicamente de tiempo disponible.

\begin{table}[H]
    \centering
    \caption{Bloques de desarrollo y criterio de cierre}
    \begin{tabular}{@{}>{\raggedright\arraybackslash}p{1.1cm}>{\raggedright\arraybackslash}p{4.6cm}>{\raggedright\arraybackslash}p{6.0cm}@{}}
        \toprule
        \textbf{N.\textsuperscript{o}} & \textbf{Objetivo} & \textbf{Criterio de cierre} \\
        \midrule
        0 & Desbloquear el módulo de backend & Contrato de interfaz entregado \\
        1 & Escaneo real en dispositivo & Prueba 1.4 superada: punto de decisión sobre la librería \\
        2 & Persistencia local & Historial y preferencias sobreviven al reinicio \\
        3 & Integración con el backend & Filtro F1 funcional de extremo a extremo con degradación controlada \\
        4 & Completar funcionalidades parciales & Reporte, Story Card e historial operativos \\
        5 & Compatibilidad con iOS & Fuera de alcance por falta de hardware \\
        6 & Pruebas y métricas & Cobertura, mediciones en hardware y pruebas de usabilidad \\
        7 & Cierre y documentación & Informe y defensa \\
        \bottomrule
    \end{tabular}
    \label{tab:bloques}
\end{table}

El bloque~1 concentra el mayor riesgo técnico del módulo, razón por la cual se
adelanta: si el escaneo no cumple los tiempos, todo lo construido encima queda
condicionado. El bloque~5 se mantiene en el plan aunque esté fuera de alcance, para
que la exclusión sea una decisión visible y no una omisión.

\subsection{El contrato de interfaz como instrumento}

La coordinación con el módulo de backend se formalizó mediante un contrato de interfaz
escrito antes de que existiera el endpoint, y conviene señalar el procedimiento con
que se redactó porque tiene valor metodológico. El contrato no se derivó de lo que el
backend podría ofrecer, sino del consumo real del frontend hacia atrás: las doce
pantallas ya implementadas sobre datos de prueba determinaban exactamente qué campos
se necesitaban, con qué tipos y en qué momento, lo que convierte a las definiciones de
tipos del frontend en una especificación ya validada contra una interfaz funcional, y
no en una lista de deseos.

El contrato declara formato JSON en UTF-8, marcas temporales en ISO~8601 e
identificación anónima de dispositivo sin cuentas de usuario, y traslada la regla de
seguridad fundamental del producto a una restricción sobre la API: ningún endpoint
puede devolver HTML para que el cliente lo renderice, ni redirigir el dispositivo
hacia el sitio analizado. A la fecha de corte de este informe permanecen ocho puntos
del contrato sin acordar, lo que se consigna como dependencia externa y no como tarea
pendiente del módulo.

\section{Estrategia de validación}

La validación se organiza en tres niveles, separados por el tipo de pregunta que cada
uno puede responder. La separación es deliberada: asignar a un nivel una pregunta que
no puede contestar es la forma más común de producir evidencia que parece sólida sin
serlo.

\subsection{Nivel 1: pruebas unitarias de la lógica de negocio}

El primer nivel cubre la capa de servicios mediante pruebas unitarias automatizadas.
A la fecha de corte el conjunto comprende 249 pruebas distribuidas en 19 suites, con
un umbral de cobertura del 70~\% en sentencias, ramas, funciones y líneas que
interrumpe la compilación si no se alcanza.

La cobertura se mide únicamente sobre la capa de servicios. Las pantallas y el sistema
de diseño quedan excluidos a propósito, no por descuido: lo que se verifica en ellos
---que un color comunique riesgo, que una jerarquía de botones induzca a la decisión
correcta--- no es una afirmación que una prueba unitaria pueda sostener, y se valida en
dispositivo y en las pruebas de usabilidad.

El aislamiento se consigue con cuatro dobles de prueba y ninguno más: sustitución
global del almacén clave-valor por una implementación en memoria con la misma
interfaz síncrona, reemplazo de la función de red del entorno por una simulada,
sustitución del cliente HTTP en las pruebas de los servicios que lo consumen y
sustitución del módulo de ubicación. En ningún momento se realizan peticiones de red
reales, de modo que el conjunto es determinista y reproducible sin conectividad ni
credenciales.

Dos pruebas del conjunto merecen mención porque validan el método y no solo el código.
La primera fija como aserción el valor del tiempo límite de la petición de análisis:
sin ella, elevarlo a 60~s dejaría pasar todas las demás pruebas sin que nadie lo
advirtiera, y un parámetro con consecuencias sobre la experiencia de uso quedaría sin
vigilancia. La segunda recorre la tabla de códigos QR de ejemplo de la documentación y
verifica que cada entrada produzca la clasificación que la tabla declara, con lo que la
documentación deja de poder divergir silenciosamente del comportamiento real.

\subsection{Nivel 2: pruebas en dispositivo físico}

El segundo nivel cubre todo lo que resulta inverificable desde una computadora:
diálogos del sistema operativo, cámara, GPS real, galería, compartición y el backend
en producción. Se compone de dos protocolos numerados que se detallan en la
sección~\ref{sec:mediciones}.

\subsection{Nivel 3: pruebas de usabilidad}

El tercer nivel evalúa lo que ningún test automatizado puede responder: si el veredicto
se entiende. Se realizará con cinco participantes y medirá tres indicadores:
comprensión del veredicto, con criterio de al menos el 90~\%; finalización del flujo
completo sin asistencia, con criterio de al menos el 85~\%; y tiempo de aprendizaje,
con criterio máximo de 60~s. El fundamento de medir la comprensión por separado del
desempeño técnico proviene de la literatura de \emph{usable security}: el modelo del
humano en el ciclo descompone el comportamiento seguro en notar la señal, entenderla y
actuar en consecuencia \cite{cranor2008framework}, y la evidencia de campo a gran
escala muestra que la tasa de continuación ante una advertencia varía según la
severidad percibida y la claridad del mensaje \cite{akhawe2013alice}. Un veredicto
correcto que no se entiende no modifica la conducta y, por lo tanto, no protege.

\section{Diseño de las mediciones en hardware}
\label{sec:mediciones}

\subsection{Prueba 1.4: tiempo de escaneo y tasa de detección}

La prueba~1.4 constituye el punto de decisión del bloque~1 y mide tres indicadores con
definiciones operativas explícitas.

La tasa de detección (TDR) se define como el número de códigos detectados dividido
entre el número de códigos presentados, multiplicado por cien, con criterio de
aceptación de al menos el 95~\%. El tiempo de escaneo (\(T_{scan}\)) se mide desde el
instante en que la pantalla del escáner queda lista hasta que la cámara entrega un
código, con criterio de aceptación de 5~s o menos en el percentil~90; excluye
deliberadamente el análisis posterior, que constituye un indicador distinto
(\(T_{F1}\)) medido contra el backend. La precisión de clasificación se define como la
proporción de códigos que abren la pantalla correcta según su tipo de contenido, con
criterio de al menos el 98~\%.

La muestra consta de 30 códigos que varían las tres condiciones que el uso real impone,
conforme al cuadro~\ref{tab:muestra}. La variación por ángulo es la que más interesa,
pues es donde típicamente se degrada la detección; la condición de oscuridad con
linterna encendida cumple además la función de verificar el control de linterna sobre
hardware real.

\begin{table}[H]
    \centering
    \caption{Composición de la muestra de la prueba 1.4}
    \begin{tabular}{@{}>{\raggedright\arraybackslash}p{2.6cm}>{\raggedright\arraybackslash}p{9.0cm}@{}}
        \toprule
        \textbf{Condición} & \textbf{Niveles} \\
        \midrule
        Tamaño & 3 pequeños ($\sim$2~cm, etiqueta o ticket); 4 medianos ($\sim$5~cm, volante o menú); 3 grandes ($\sim$15~cm o más, cartel a distancia) \\
        Ángulo & 4 de frente; 3 inclinados $\sim$30\textdegree; 3 inclinados $\sim$45\textdegree \\
        Iluminación & 4 con luz de día; 3 en penumbra; 3 en oscuridad con linterna encendida \\
        Contenido & Al menos 20 con URL; 5 de credenciales Wi-Fi; 5 de texto plano o contacto \\
        \bottomrule
    \end{tabular}
    \label{tab:muestra}
\end{table}

Los últimos diez códigos, de contenido no navegable, son los que validan la métrica de
clasificación: lo que se comprueba en ellos es la lectura por cámara, no el análisis
sintáctico del contenido, que ya se verifica en el nivel~1.

La prueba se ejecuta sobre un teléfono Android de gama media y no sobre el mejor
dispositivo disponible, dado que el criterio de tiempo de arranque está definido para
gama media y conviene medir sobre hardware representativo. La instrumentación reside en
la propia aplicación: cada lectura registra su duración y el resumen calcula la mediana,
el percentil~90, el máximo, la tasa de detección y el cumplimiento de ambos criterios.
El percentil se obtiene por interpolación lineal. Si un código no se lee en 10~s se
registra como fallo y se continúa con el siguiente; ese fallo cuenta contra la tasa de
detección.

Se eligió el percentil~90 y no el promedio ni el máximo por una razón que conviene
dejar asentada: una única lectura lenta no debe descalificar la aplicación, pero
tampoco debe poder esconderse en un promedio. Con nueve lecturas de 500~ms y una de
20~s el percentil~90 sigue cumpliendo el criterio; hacen falta dos lecturas lentas para
desplazarlo.

Corresponde declarar una limitación del instrumento. Las lecturas exitosas se
registran solas, pero los fallos deben anotarse manualmente, porque la aplicación no
tiene forma de saber que hubo un código frente a la cámara que nunca llegó a
decodificar. Si los fallos no se anotan, la tasa de detección resulta del 100~\% y el
dato es falso. Es el único valor que el instrumento no puede deducir y por eso el
protocolo lo advierte de forma destacada. Adicionalmente, las muestras residen en
memoria y se pierden al recargar la aplicación, de modo que los tiempos se anotan
conforme se producen.

\subsection{Protocolo de verificación en dispositivo}

El segundo protocolo es una lista de verificación de aproximadamente 50 minutos,
pensada para ejecutarse en una sola sesión, organizada en cinco bloques con resultado
esperado por paso. Cubre el arranque en frío y el onboarding completo; el análisis real
contra el backend, incluyendo los tres veredictos, el comportamiento ante una URL no
vista antes y el aviso de espera prolongada; la geocodificación inversa, con criterio
de al menos el 85~\% de nombres de lugar legibles, que debe ejecutarse al aire libre;
la generación de la Story Card, con criterio de al menos el 90~\%; el camino sin
conexión; y el comportamiento del teclado y las áreas seguras de la pantalla.

Dos pasos del protocolo están marcados como no omisibles. El primero es de seguridad:
consiste en intentar escanear con un segundo teléfono el código QR que la Story Card
muestra cubierto; si resulta legible, se trata de un fallo grave, y es la razón por la
que el código se cubre con una capa opaca y no con un difuminado, pues un código
borroso sigue siendo escaneable. El segundo verifica que la preferencia de ubicación
sobreviva al reinicio, regresión que se detectó y corrigió durante el desarrollo.

La instrucción transversal del protocolo es registrar qué paso falló y qué se observó,
en lugar de marcar únicamente aprobado o reprobado: un paso fallido acompañado del dato
de lo ocurrido vale más que diez marcados en verde.

\section{Latencia del backend y política de degradación}

El tiempo límite de la petición de análisis se fijó en 25~s a partir de mediciones
sucesivas contra el backend real, que se consignan porque explican una decisión que de
otro modo parecería contradecir el criterio de 5~s del indicador \(T_{F1}\). En la
primera medición, la llamada inicial tomó 5,5~s por arranque en frío de la función sin
servidor y entre 1,6 y 3,1~s en las siguientes. Diez días después, una URL no vista
antes por el backend tomaba entre 14,2 y 14,5~s, frente a 0,3~s si ya estaba en caché;
el salto coincide con la incorporación del filtro de vista previa visual, que abre un
navegador sin interfaz y cuesta alrededor de 10~s. Una medición posterior confirmó
15,3~s en una URL nueva.

Con un límite de 12~s, todo escaneo de un sitio nuevo se interrumpía solo y degradaba
a veredicto incierto, es decir, precisamente el caso normal, porque un código QR de
interés apunta casi siempre a algo que el backend no ha visto antes. El límite de 25~s
deja cerca de 11~s de margen sobre el peor caso medido, y la pantalla de carga avisa al
usuario a los 8~s que la espera se está prolongando.

El argumento metodológico que sostiene esta decisión conviene explicitarlo: el tiempo
límite no es la métrica. El indicador \(T_{F1}\) mide lo que tarda el backend en
responder, no el momento en que la aplicación se rinde. Recortar el límite no mejora el
número medido; solo convierte análisis buenos en inciertos. Si la latencia del backend
disminuye, el límite deberá medirse y reducirse de nuevo, pues 25~s de espera es mucho
pedir a una persona que sostiene el teléfono en alto.

\subsection{Taxonomía de errores y degradación}

El cliente HTTP traduce toda condición de fallo a un vocabulario cerrado de once
códigos, de modo que la interfaz decida siempre según el código y nunca interpretando
el texto del mensaje. La traducción absorbe dos formas de cuerpo de respuesta
incompatibles ---la acordada en el contrato y la que el endpoint devuelve en la
práctica--- y normaliza los nombres alternativos que el backend emplea. El cuadro
\ref{tab:errores} resume el criterio de reintento, que la interfaz usa para decidir si
ofrece volver a intentar.

\begin{table}[H]
    \centering
    \caption{Clasificación de errores según posibilidad de reintento}
    \begin{tabular}{@{}>{\raggedright\arraybackslash}p{3.0cm}>{\raggedright\arraybackslash}p{8.6cm}@{}}
        \toprule
        \textbf{Categoría} & \textbf{Códigos y tratamiento} \\
        \midrule
        Reintentable & Tiempo agotado del filtro, filtro no disponible, límite de tasa excedido, error interno, fallo de red y tiempo agotado de la petición \\
        No reintentable & URL inválida, carga inválida y análisis no encontrado: el fallo está en la petición, de modo que repetirla produciría el mismo resultado \\
        Cancelación & Se distingue del tiempo agotado aunque el mecanismo subyacente sea el mismo, pues abandonar la pantalla no es un error y no debe reportarse como tal \\
        \bottomrule
    \end{tabular}
    \label{tab:errores}
\end{table}

La política de degradación distingue dos casos. Una URL inválida se repropaga al
llamador, porque no constituye un fallo de análisis sino contenido no analizable, y la
pantalla lo trata como tal dirigiendo al usuario a la vista de contenido. Cualquier otro
error degrada a un resultado honesto: veredicto incierto con aviso visible de que no se
completaron todas las verificaciones. En ningún caso un fallo produce un veredicto
seguro, criterio que se deriva de la regla de que la falta de información nunca se
comunica en verde.

\section{Determinación del veredicto}

El veredicto se resuelve con una precedencia de dos fuentes. Si el backend concluyó, su
veredicto manda; un valor pendiente no constituye conclusión, mientras que un incierto
explícito sí se respeta. Solo si el backend no concluyó se deriva el veredicto del
puntaje numérico.

Esta derivación existe por una razón circunstancial que conviene documentar: el backend
calcula correctamente el puntaje y sus criterios, pero su campo de veredicto final
responde pendiente a toda consulta, incluidos dominios legítimos, porque el umbral que
convierte puntaje en veredicto no está conectado de su lado. Sin la derivación, todo
análisis terminaría en incierto y la aplicación no podría comunicar lo único que debe
comunicar. No se trata de un veredicto inventado: el puntaje y los criterios que lo
componen provienen del backend; lo único que aporta la aplicación son los cortes. El
diseño se desmantela por sí solo, pues en cuanto el backend devuelva veredictos reales
la primera fuente tiene precedencia y la derivación deja de ejecutarse sin cambiar una
línea.

Los cortes están alineados con los pesos que el propio backend asigna a cada criterio
---dominio de nivel superior de riesgo, 25; host que es una dirección IP literal, 30;
ausencia de HTTPS, 15; palabra clave sospechosa, 15--- según el cuadro
\ref{tab:bandas}.

\begin{table}[H]
    \centering
    \caption{Bandas de puntaje y su fundamento}
    \begin{tabular}{@{}>{\raggedright\arraybackslash}p{1.8cm}>{\raggedright\arraybackslash}p{2.2cm}>{\raggedright\arraybackslash}p{7.6cm}@{}}
        \toprule
        \textbf{Puntaje} & \textbf{Veredicto} & \textbf{Fundamento} \\
        \midrule
        0--29 & Seguro & Un solo criterio menor no basta. Marcar de riesgo todo sitio sin HTTPS produciría tantos falsos positivos que la aplicación perdería credibilidad \\
        30--59 & Sospechoso & 30 es exactamente el peso de un host que es una IP literal, y también la suma de dos criterios menores. Merece advertencia, no condena \\
        60--100 & Malicioso & Tres o más señales simultáneas, por ejemplo IP más ausencia de HTTPS más la palabra \emph{login}: el caso clásico de phishing \\
        \bottomrule
    \end{tabular}
    \label{tab:bandas}
\end{table}

\subsection{Validación de rango}

El puntaje se valida en dos puntos independientes del flujo, uno para la derivación del
veredicto y otro para el medidor visual de riesgo. La duplicación es intencional y
ambas validaciones fallan hacia el lado seguro.

Un valor ausente o que no es un número finito no produce veredicto optimista, sino
ninguno, y el llamador lo resuelve como incierto. Un valor fuera del rango de 0 a 100
también se rechaza, y por la misma razón de fondo: un $-20$ o un $150$ no significan
``muy seguro'' y ``muy malicioso'', sino que el backend cambió de escala y que los
cortes ya no significan lo que se supone. Interpretarlos con la tabla anterior es
exactamente el mecanismo por el cual se produce un veredicto equivocado con apariencia
de correcto, dado que un $-20$ daría seguro, y un falso seguro es el peor resultado
posible del producto. Ante una escala que no se comprende, la aplicación informa
incierto, que es la afirmación verdadera.

\section{Métricas e instrumentos de evaluación}

El cuadro~\ref{tab:metricas} asocia cada métrica del anteproyecto con el instrumento
que la medirá. La columna de instrumento es la aportación metodológica del cuadro: sin
ella, una métrica puede enunciarse sin que exista procedimiento capaz de producirla.

\begin{table}[H]
    \centering
    \caption{Métricas del proyecto y su instrumento de medición}
    \begin{tabular}{@{}>{\raggedright\arraybackslash}p{4.3cm}>{\raggedright\arraybackslash}p{2.3cm}>{\raggedright\arraybackslash}p{5.0cm}@{}}
        \toprule
        \textbf{Métrica} & \textbf{Criterio} & \textbf{Instrumento} \\
        \midrule
        Clasificación URL / no-URL & $\geq$ 98~\% & Pruebas unitarias y prueba 1.4 \\
        Cobertura de pruebas unitarias & $\geq$ 70~\% & Informe de cobertura automatizado \\
        Persistencia offline del historial & 100~\% & Pruebas unitarias y protocolo en dispositivo \\
        Reportes con carga completa & 100~\% & Pruebas unitarias \\
        \(T_{scan}\) (percentil 90) & $\leq$ 5~s & Prueba 1.4 en dispositivo \\
        Tasa de detección (TDR) & $\geq$ 95~\% & Prueba 1.4 en dispositivo \\
        \(T_{F1}\) (percentil 90) & $\leq$ 5~s & Medición contra el backend \\
        \(T_{total}\) & $\leq$ 35~s & Medición contra el backend \\
        Detección de maliciosos & $\geq$ 90~\% & Muestra de URLs verificadas, requiere backend \\
        Falsos positivos & $\leq$ 5~\% & Muestra de dominios legítimos, requiere backend \\
        Geocodificación legible & $\geq$ 85~\% & Protocolo en dispositivo, bloque B \\
        Story Card generada & $\geq$ 90~\% & Protocolo en dispositivo, bloque C \\
        Comprensión del veredicto & $\geq$ 90~\% & Prueba de usabilidad \\
        Finalización del flujo & $\geq$ 85~\% & Prueba de usabilidad \\
        \(T_{learn}\) & $\leq$ 60~s & Prueba de usabilidad \\
        \(T_{start}\) & $\leq$ 3~s & Medición en dispositivo de gama media \\
        Compatibilidad Android $\geq$ 10 & Sin errores críticos & Compilación nativa y protocolo en dispositivo \\
        \bottomrule
    \end{tabular}
    \label{tab:metricas}
\end{table}

Cuatro de las diecisiete métricas quedan cubiertas por el nivel de pruebas
automatizadas. Las trece restantes dependen de hardware, del avance del módulo de
backend o de la disponibilidad de participantes, y así se declaran.

Un caso exige precisión particular. La métrica de detección de maliciosos no es medible
en el estado actual del sistema, y no por una limitación del instrumento: el backend no
devuelve el veredicto malicioso en ninguna consulta, de modo que la medición arrojaría
0~\% por construcción y no por desempeño. La métrica queda condicionada a que el
módulo de backend conecte su umbral de decisión, y su valor no debe interpretarse
hasta entonces. Las métricas de clasificación que sí dependen del frontend se evalúan
con las definiciones estándar de precisión y exhaustividad para clasificadores
\cite{sokolova2009systematic}.

\section{Consideraciones éticas y de privacidad}

La minimización de datos se trata aquí como decisión de método y no como añadido, pues
determina qué puede medirse y cómo.

La regla de mayor jerarquía del producto es que la aplicación nunca visita la URL
analizada. Esta regla domina sobre cualquier otra decisión de producto y tiene
consecuencia metodológica directa: el análisis opera sobre la cadena de texto, de modo
que ninguna medición de este trabajo implica contacto con la infraestructura del
atacante, ni expone la dirección IP del dispositivo, ni revela que un usuario concreto
está por acceder a una URL concreta. La vista previa visual, cuando el backend la
provee, se presenta como imagen estática y nunca como contenido navegable.

De esa regla se derivan las demás decisiones. La geocodificación inversa se resuelve en
el dispositivo y no en el servidor, de modo que las coordenadas no salen del teléfono.
Los reportes de códigos maliciosos son anónimos: no incluyen identificador de usuario
ni coordenadas, y los campos opcionales solo se transmiten cuando contienen datos. El
historial reside exclusivamente en el teléfono, que es la fuente de verdad, lo que
además es la condición que hace posible la métrica de recuperación sin conexión. Las
contraseñas de redes Wi-Fi codificadas en un QR se enmascaran en pantalla. El código
QR que aparece en la Story Card se cubre con una capa opaca y no con un difuminado,
porque un código borroso sigue siendo escaneable y compartir la tarjeta propagaría el
ataque.

El marco de verificación de seguridad de la aplicación es el estándar de verificación
de seguridad de aplicaciones móviles de OWASP en su versión 2.1
\cite{owasp2024masvs}, del que aplican tres categorías al prototipo: almacenamiento,
para el historial local; red, para la comunicación con el backend; y plataforma, para
el manejo de permisos de cámara y ubicación. Se eligió por ofrecer criterios
verificables en lugar de buenas prácticas generales.

Para las pruebas de usabilidad se recabará consentimiento informado, se explicará el
propósito del estudio y el carácter voluntario de la participación, y los datos se
tratarán de forma anónima y agregada. Los participantes escanearán únicamente códigos
preparados para la prueba, nunca códigos maliciosos reales en entornos de uso
cotidiano.

\section{Limitaciones del método}

Corresponde declarar las limitaciones del diseño, con la misma precisión con que se
declararon las de la revisión de antecedentes.

La cobertura automatizada alcanza la lógica de negocio pero no las pantallas: no hay
pruebas de componentes ni pruebas de extremo a extremo automatizadas, de modo que el
comportamiento de la interfaz se verifica mediante protocolos manuales en dispositivo,
cuya reproducibilidad depende de quien los ejecute. Tampoco hay analizador estático de
estilo configurado: la puerta de calidad del proyecto son la verificación de tipos en
modo estricto y el conjunto de pruebas.

La plataforma iOS queda fuera de alcance por falta de hardware, conforme a lo
declarado en el capítulo de alcance, de modo que las conclusiones sobre desempeño del
escaneo se limitan a Android.

La muestra de las pruebas de usabilidad es de cinco participantes, insuficiente para
inferencia estadística. Sus resultados se interpretarán como indicios de problemas de
comprensión y no como estimaciones poblacionales, y el estudio de usabilidad con
muestra representativa ya figura entre las exclusiones explícitas del trabajo.

Las mediciones de latencia del backend corresponden a un endpoint en ambiente de
desarrollo con recursos del nivel gratuito, de modo que los tiempos observados incluyen
arranques en frío que un despliegue en producción podría no presentar. El filtro de
vista previa visual se comporta de forma intermitente, por lo que no se cuenta con él
para las mediciones.

Finalmente, a la fecha de corte de este informe la aplicación no se ha ejecutado en un
teléfono físico: no hubo dispositivo ni emulador disponible en el entorno de
desarrollo. Lo verificado hasta ahora es que la verificación de tipos está limpia, que
las 249 pruebas pasan, que el paquete JavaScript de la aplicación completa se exporta
sin errores para Android ---lo que garantiza ausencia de importaciones rotas, rutas mal
declaradas y errores de sintaxis en las doce pantallas, pero nada de lo que solo se
observa ejecutando--- y que los tres veredictos se producen correctamente al procesar
respuestas en vivo del backend real. Todo lo que requiere hardware permanece pendiente
de los protocolos descritos en la sección~\ref{sec:mediciones}, y así se consigna en
lugar de presentarse como verificado.
